Kita tulis

\mavergleichskette
{\vergleichskette
{ w
}
{ = }{ a+b { \mathrm i}
}
{  }{
}
{    }{
}
{   }{
}
}
{}{}{.}
Maka

\mavergleichskettealign
{\vergleichskettealign
{  \psi ( a+b { \mathrm i}   )
}
{ =} { { \frac{   \left(  a+b { \mathrm i} +1  \right) { \mathrm i}      }{  1-  a-b { \mathrm i}      } }
}
{ =} { { \frac{     a { \mathrm i} +{ \mathrm i}  -b  }{  1-  a-b { \mathrm i}      } }
}
{ =} { { \frac{     \left(  a { \mathrm i} +{ \mathrm i}  -b  \right) \left(  1-  a+b { \mathrm i}  \right)  }{  \left(  1-  a-b { \mathrm i}  \right) \left(  1-  a+b { \mathrm i}  \right)      } }
}
{ =} { { \frac{     \left(  a+1 \right) \left(  1-a  \right) { \mathrm i}  -b^2 { \mathrm i} -b \left(  1-  a \right) -b\left(  a+1  \right)  }{  \left(  1-  a \right)^2 +b^2  } }
}
}
{
\vergleichskettefortsetzungalign
{ =} { { \frac{     \left(  1-a ^2-b^2  \right) { \mathrm i}   -2 b  }{  \left(  1-  a \right)^2 +b^2  } }
}
{ } {
}
{ } {}
{ } {}
}
{}{.}
Jadi, dalam koordinat real,

\mavergleichskettedisp
{\vergleichskette
{ \psi \begin{pmatrix}  a  \\b   \end{pmatrix}
}
{ =} { { \frac{   1  }{  -2a+1+a^2+b^2 } } \begin{pmatrix}  -2b \\ 1-a^2-b^2   \end{pmatrix}
}
{ } {
}
{ } {
}
{ } {
}
}
{}{}{.}
